\documentclass[12pt]{article}

\usepackage{graphicx}
\usepackage{amsfonts}
\usepackage{amsthm}
\usepackage{amsmath}
\usepackage{amssymb}
\usepackage{scribe}
\usepackage{braket}
\usepackage{xcolor}
\usepackage{nicematrix}

\usepackage[margin=2in]{geometry}
\pagestyle{empty}

\DeclareMathOperator{\rk}{rk}

\begin{document}

\makeheader{Kevin Sheng}{Midterm}

have you ever communicated bits with your life on the line?

\section*{Problem 1}

thank god fooling set worked for this one

I propose that sending the $O(n^2)$ adjacency matrix straight up is the best we can do.

Consider the set of equivalence classes $I$ on all the graphs,
where two graphs are equivalent if they're isomorphic.
I'll use $[G]$ to denote the equivalence class of a graph $G$.

Notice that $\{(G, G) \mid [G] \in I\}$ makes a fooling set
simply by virtue of how we defined $I$ in the first place.
Any graph is always isomorphic to itself, and if two graphs are in
different classes then they can't be isomorphic to each other.

Now let us lower bound the size of this set.
In total, there's $2^{\binom{n}{2}}$ graphs, and each graph
can be isomorphic to at most $n!$ other graphs.
To see this, notice that an isomorphism is essentially a relabelling/permutation
of the original nodes, of which there are $n!$ of.

Thus, $|I| \ge \frac{2^{\binom{n}{2}}}{n!}$ and we have
\begin{align*}
    \mathcal{D}(\text{graph iso})
     & \ge \log |I|                       \\
     & = \log 2^{\binom{n}{2}} - \log n!  \\
     & = \binom{n}{2} \log 2 - \log n!    \\
     & \ge \binom{n}{2} \log 2 - n \log n \\
     & = \mathcal{O}(n^2)\quad\square
\end{align*}

\newpage

\section*{Problem 2}

no shot fooling sets work here too

WLOG let Alice and Bob's numbers be less than $2^n$ so we can
consider $f: \{0, 1\}^n \times \{0, 1\}^n \to \{0, 1\}$.
Any bits larger than that won't affect the result of the $n$th bit anyways.

Notice that $x$ has a $0$ in the $n$th position if $0 \le x < 2^{n-1}$ and a $1$ otherwise.
Thus, let's consider the following set:
\[S=\left\{\left(x, \left\lfloor\frac{2^{n-1}-1}{x}\right\rfloor\right)
\middle| 1 \le x(x+1) < 2^{n-1}\right\}\]
For example, if $2^{n-1}=16$,
\[S=\{(1, 15), (2, 7), (3, 4)\}\]
I'm pretty sure this set has around $2^{\frac{n-1}{2}}$ elements.

By construction, the product of every pair in this set evaluates to $0$.
Now take $(a, b)$ and $(c, d)$, with $a < c$.
We need $bc \ge 2^{n-1}$ for each crosspair to be nonconstant.
Producing the following lower bound on the product:
\begin{align*}
    \left\lfloor\frac{2^{n-1}-1}{a}\right\rfloor \cdot (a+(c-a))
    &\ge \left\lfloor\frac{2^{n-1}-1}{a}\right\rfloor \cdot (a+1) \\
    &= \left\lfloor\frac{2^{n-1}-1}{a}\right\rfloor \cdot a + \left\lfloor\frac{2^{n-1}-1}{a}\right\rfloor \\
    &\ge 2^{n-1}-1-a + \left\lfloor\frac{2^{n-1}-1}{a}\right\rfloor
\end{align*}
means we need $\left\lfloor\frac{2^{n-1}-1}{a}\right\rfloor \ge a+1$.
However, this condition was already satisfied by our construction, as $a(a+1) < 2^{n-1}$.

Since any pair gives a result of $0$ and any crosspair has a result of $1$,
this is a $0$-fooling set and we have a lower bound of on
the communication complexity of $\mathcal{O}(n)$.

\newpage

\section*{Problem 3}

I believe that any $\phi$ can be converted into CNF where for each $z_i$,
its appearance in the formula is either all $z_i$ or all $\lnot z_i$.
So something like this:
\[(z_1 \lor \lnot z_2) \land (\lnot z_2 \lor z_3)\]
would be fine, but not
\[(z_1 \lor \lnot z_2) \land (z_2 \lor z_3)\]
since both $z_2$ and $\lnot z_2$ appear in the formula.

The intuition for this can be seen by induction on the DNF forms.
If there's only one leaf, then we only have $z_1$ and $\lnot z_1$.
Now we assume that any formula involving $z_1$ through $z_n$
can be written like
\[\bigvee_{i=1}^{N} \bigwedge_{j \in S_i} e_j\]
where $N$ is whatever number of summands, $S_i$ is all the coefficients
of the $i$th summand, and $e_j$ is either always $z_j$ or $\lnot z_j$.

Now consider a formula on $e_1, e_2, \cdots, e_{n+1}$.
Since boolean operations are associative, we can turn any formula tree into
a binary tree and thus WLOG assume that $e_{n+1}$ is immediately combined
through AND or OR with some $e_x$.
This gives us a formula on $e_1, e_2, \cdots, e_n, e_{n+1}e_x$ or a formula
on $e_1, e_2, \cdots, e_n, (e_{n+1}+e_x)$.
Either way, we can use distributivity laws on the formula given
by our inductive hypothesis to get something in the form that we want.
Nowhere in the process do we need both $e_{n+1}$ and $\lnot e_{n+1}$ at the same time.

By laws of boolean algebra, I'm pretty sure we can then turn this into the CNF form of
\[\bigwedge_{i=1}^{M} \bigvee_{j \in S_i} e_j\]

Now consider the characteristic matrix $M$:
\[M = \left[\prod_{i=1}^{M} \bigvee_{j \in S_i} e_j\right]_{(x_1, \cdots, x_n),(y_1, \cdots, y_n)}\]
where I've converted AND to $\prod$ since in $\mathbb{F}_2$ the two are the same thing.

Let's narrow down to one $M_i=\left[\bigvee_{j \in S_i} e_j\right]_{x_{S_i}, y_{S_i}}$.
I propose that the structure of this matrix can be arranged into something like
\[\begin{bmatrix}
        0 & 1 & 1 & 1 \\
        1 & 0 & 1 & 1 \\
        1 & 1 & 0 & 1 \\
        1 & 1 & 1 & 0
    \end{bmatrix}\]
that has rank $2^{|S_i|}-1$ if $|S_i| > 1$ and full rank otherwise.

To see this, consider a fixed $x$.
The only way for this matrix to have a $0$ at the associated $y$ is if $e_j=0\ \forall j \in S_j$.
Whether $e_j=x_j \oplus y_j$ or $e_j=\lnot (x_j \oplus y_j)$, there is only
one value of $y_j$ that satisfies all of these constraints.
Continuing this train of thought, we see that if $x_1 \ne x_2$, then their
associated $y_1$ and $y_2$ that give a zero entry must be different as well.
To prove this, consider the bit $k \in S_i$ at which $x_1$ and $x_2$ differ.
WLOG let $(x_1)_k = 1$ and $(x_2)_k=0$.
If $e_k=x_k \oplus y_k$, then this forces $(y_1)_k=1$ and $(y_2)_k=0$, forcing $y_1 \ne y_2$.
By symmetry, the same holds if $e_k=\lnot (x_k \oplus y_k)$.

If all $S_i$s were disjoint, we're done.
We could take the massive product out to get
\[M = \bigotimes_{i=1}^{N} \left[\bigvee_{j \in S_i} e_j\right]_{x_{S_i}, y_{S_i}}\]
and use $\rk A \otimes B = \rk A \cdot \rk B$.
Since each increase in $n$ by $1$ multiplies the rank by at least $1.5$,
$\rk M \ge 1.5^n$ and $\mathcal{D}(\phi) = \Omega(n)$ by extension.

However, sometimes the $S_i$s aren't disjoint.
In this case I'm lowkey at a loss for logic.
The idea is to first insert "dummy indices", where for each $e_j$
I split it into multiple $e_{ij}$s where $i$ is the specific $S_i$ that this $e_j$ belongs to.
This temporarily makes all $S_i$s disjoint, so I can still do the matrix product reduction.

Now consider an $e_k$ that's in both $S_i$ and $S_j$.
(Hopefully my sudden new overload for $j$ isn't too confusing.)
If we organize the matrix with respect to $(e_{ik}, e_{jk})$ as our
central values, we get something like this:
\[M'=\begin{bNiceArray}{cc|cc}[first-row,first-col, margin]
        & (0, 0)         & (1, 1) & (0, 1)         & (1, 0) \\
        (0, 0) & \Block{2-2}{A} &        & \Block{2-2}{B} &        \\
        (1, 1) &                &        &                &        \\
        \hline
        (0, 1) & \Block{2-2}{B} &        & \Block{2-2}{A}          \\
        (1, 0) &                &        &                &        \\
    \end{bNiceArray}\]
Since $(0, 0) \oplus (0, 0) = (0, 1) \oplus (0, 1)$ and so on and so forth,
the second quadrant is actually equivalent to the fourth quadrant.
Same can be said for the first and third quadrant.

In our actual $M$, we only want to select $A$, but I'm not sure how this will affect the rank.
I'm pretty sure it's just going to halve the rank in the worst case, but
I have no clue how to prove that.

If doing this does indeed only halve the rank, though, then we can reconstruct
$M$ by continually parsing it from the remnants of wherever $M'$ was and then
see that the log of the rank is still on the order of $\Omega(n)$.

\newpage

\section*{Problem 4}

The upper bound is $\mathcal{O}(n)$ as always.
It remains to prove a lower bound on $\mathcal{N}(f)$.

The idea is to use the RS bound.
Let's first try and get an estimate on the number of orthogonal pairs,
or the number of $1$s in the characteristic matrix.

I'm pretty sure we can do this by calculating the chance of this event:
\[P\left(\sum_{i=1}^{n} x_iy_i \equiv 0 \pmod{18181}\right)\]
Each $x_i$ and $y_i$ is chosen independently,
so this is equivalent to $P(X \equiv 0 \pmod{18181})$,
where $X \sim \text{Binom}\left(n, \frac{1}{4}\right)$.
By symmetry, as $n \to \infty$ we get that $P(X \equiv 0 \pmod{18181})=\frac{1}{18181}$,
since there's no reason to favor a specific value mod $18181$ as opposed
to any other remainder.

Thus, the number of ones is pretty close to $\frac{4^n}{18181}$.
Let $\mu$ be the uniform distribution over these cells, with each cell
carrying a probability of $\frac{18181}{4^n}$.

Now we can carry over the reasoning we did for $\mathbb{F}_2$.
Any maximal rectangle with only $1$s must consist of the product
of two subspaces, as if $\braket{a, b}=0$ for all $a \in A$ then the same
must hold for all $a \in \text{span}(A)$ as well.
Similar reasoning holds for $B$ and $\text{span}(B)$.

Since $\text{span}(A)$ and $\text{span}(B)$ are two orthogonal subspaces,
the sum of their dimensions cannot exceed $n$ and so
\[|A \times B| \le 2^{\dim \text{span}(A)}2^{\dim \text{span}(B)} \le 2^n\]
This captures at most $2^n \cdot \frac{18181}{4^n}=\frac{18181}{2^n}$ of the probability,
so $\text{RS}(f) \le \frac{18181}{2^n}$ and by extension
$\mathcal{N}(f) \ge \log \frac{1}{\text{RS}(f)} = \Theta(n)$. $\square$

\newpage

\section*{Problem 5}

There's some facts about prime number density that I'm going to use here.
I got these while doing competitive programming:
\begin{itemize}
    \item Asymptotically, the number of prime numbers below $n$ approaches $\frac{n}{\ln n}$.
    \item The maximum number of prime factors $n$ can have approaches $\frac{\log n}{\log \log n}$.
\end{itemize}

Notice that Alice and Bob can reduce their inputs to sets of prime factors that
divide their numbers instead.
This loses no information and actually slightly reduces the number of inputs
under consideration, as their numbers are coprime iff their sets of prime
factors are disjoint.

There's around $n'=\frac{2^n}{\ln 2^n}$ primes below $2^n$, and the largest
amount of prime factors either Alice of Bob can have is
$k'=\frac{\log 2^n}{\log \log 2^n} \approx \frac{n}{\log n}$.
Thus, we can use the same randomized cover construction method used for $DISJ_{n,k}$
in class to see that $\log \mathcal{C}(f)$ is around the order of
\[k'+\log \log n'=\frac{n}{\log n}+\log \log \frac{2^n}{\ln 2^n} = \Theta\left(\frac{n}{\log n}\right)\]
This serves as both a tight upper and lower bound, as $\mathcal{N}(f)$ is bounded
in both directions and below by $\log \mathcal{C}(f)$. $\square$

\newpage

\section*{Problem 6}

That $\mathcal{O}(n^2)$ is an upper bound arises from that any subspace can be represented
by a set of at most $n$ basis vectors, each with exactly $n$ bits.
Even if the problem asked for deterministic CC, we'd still have this as an upper bound.

For a lower bound, a central issue is how many subspaces over $\mathbb{F}_2^n$ there are.
To do this, let's first fix a target dimension $k$ and try to see
how many ordered sets of $k$ linearly independent vectors there are.
The first element in this set can be any nonzero vector, so we have $2^n - 1$ choices there.
The second element can be any nonzero vector that isn't the first one,
so we would have $2^n-2$ choices.
In general, when choosing the $i$th vector ($i$ starting from $0$) we'd have $2^n-2^i$ choices,
since the previous $i-1$ vectors take away $2^i$ of our options.
Thus, in total we have $\prod_{i=0}^{k-1} 2^n-2^i$ ordered bases.

However, this overcounts, as a single subspace can have many different bases.
To count how many different bases a subspace of dimension $k$ has,
we use the same process on a subspace of $2^k$ vectors to see that
a subspace of dimension $k$ has $\prod_{i=0}^{k-1} 2^k-2^i$ ordered bases.
Dividing our previous result by this gives us that $\mathbb{F}_2^n$ has
\[D_n(k)=\prod_{i=0}^{k-1} \frac{2^n-2^i}{2^k-2^i}\]
bases of dimension $k$.
Also let $P(n)=\sum_{i=0}^{n} D_n(i)$, or the total number of subspaces of $\mathbb{F}_2^n$.

To lower bound $D_n(k)$, we see that the lowest numerator is $2^n-2^{k-1}$ and
the highest denominator is $2^k-1$, which we'll treat as $2^k$ for convenience.
This leads to that
\[D_n(k) \ge \left(\frac{2^n-2^{k-1}}{2^k}\right)^k \ge \left(2^{n-k-1}\right)^k=2^{kn-k^2-k}\]
and $P(n) \ge D_n\left(\frac{n}{2}\right) \ge 2^{\frac{n^2}{4}-\frac{n}{2}}$ by extension.

OTOH, we can get an upper bound for $D_n(k)$ by seeing that the largest
numerator is $2^n-2^0 \approx 2^n$, and the smallest denominator is $2^k-2^{k-1}=2^{k-1}$.
This bounds $D_n(k)$ at the top by
\[D_n(k) \le \left(\frac{2^n}{2^{k-1}}\right)^k = 2^{nk-k^2+k}\]
Across all $k$, the exponent is maximized when $k=\frac{n}{2}$, so we get that
$P(n) \le (n+1) \cdot 2^{\frac{n^2}{4}+\frac{n}{2}}$.

Assuming a uniform distribution on $\mu$, we want to show that
\[\frac{P(n-a)P(a)}{\sum_{k=0}^{n} D_n(k)P(n-k)} \approx \frac{1}{2^{n^2}}\]
since a maximal rectangle consists of all subspaces of some subspace $S$ of dimension $a$
crossed with all subspaces of $S^\perp$ (surely :clueless:).
The denominator comes from iterating over all subspaces of dimension $k$
and then taking all subspaces of the remaining $n-k$ dimensions.

With our previous UB on $P(n)$, I'm pretty sure the numerator is maximized
when $a=0$, since that's what maximizes the exponent which dominates.
We can also LB our denominator by just taking $k=\frac{n}{2}$ to see that
\[\sum_{k=0}^{n} D_n(k)P(n-k) \ge D_n\left(\frac{n}{2}\right)D_{n/2}\left(\frac{n}{4}\right) \ge 2^{\frac{5n^2}{16}-\frac{3n}{4}}\]
Putting $P(n-0)P(n) = P(n)$ on top gives
\[\text{RS}(f) \le \frac{(n+1) \cdot 2^{\frac{n^2}{4}+\frac{n}{2}}}{2^{\frac{5n^2}{16}-\frac{3n}{4}}} = (n+1)2^{-\frac{n^2}{16}+\frac{3n}{4}}\]
and taking the log of $\frac{1}{\text{RS}(f)}$ gives $\mathcal{N}(f) \ge \Omega\left(n^2\right)$. $\square$

(proof built on genuine hopes and dreams bro gg)

\newpage

\section*{Problem 7}

We need to show that if we want a randomized protocol built out of
deterministic protocols of cost $c$, we need at least $2^{\log n/c} = \frac{n}{c}$
of them to get a correctness rate of $\frac{2}{3}$.

Each protocol is a tree of depth $c$ with $2^c$ leaves.
For each input, the randomized protocol basically assigns it a leaf in each of the protocols.
The issue now is to determine what necessitates two pairs of inputs
being assigned to different sets of leaves.
For the statement to be proven, our randomized protocol needs
at least $2^n$ distinct leaf assignments across all pairs of inputs.

Assuming this is true, if we have $p$ protocols,
there would be $(2^c)^p$ possible distinct leaf assignments.
We need this number to be at least $2^n$, so setting $p=\frac{n}{c}$ necessitates
at least $\log_2 \frac{n}{c}$ random bits.

I genuinely have no idea how to rigorously prove the lower bound of $2^n$ leaf assignments, though.
If we limit ourselves to hashing-based protocols, some kind of vision can be seen by
fixing Bob's input and considering the full range of Alice's inputs.
For any two inputs $a_1, a_2 \in \{0, 1\}^n$, there must exist at least
one hash function defined by our random bits that send $a_1$ and $a_2$
to different elements in the hashed domain, which then result in Alice
transmitting different bits to Bob.
This then results in at least one sub-deterministic protocol that assigns
different leaves to $(a_1, b)$ and $(a_2, b)$.
Alice has $2^n$ possible inputs, so the protocols need at least $2^n$
distinct leaf assignments.

\newpage

\section*{Problem 8}

Since the functions are chosen uniformly at random, maybe it'll help if we just fix some input
for Alice and consider the full range of Bob's inputs.
She knows 100\% where she is, so from Alice's POV the matrix
at this point is just a single row.

This might be a dead end, though.
\begin{enumerate}
    \item Say we split Bob's inputs into $\sqrt{2^n}$ blocks, each of size $\sqrt{2^n}$.
          If we choose a random block and have Bob tell Alice whether that was it,
          we have a $\frac{1}{\sqrt{2^n}}$ chance of being right at the cost of one bit.
          However, after that it's just cooked since now we have no way of knowing
          where Bob actually was in the block.

    \item If Alice tried to directly guess where Bob was, the problem reduces
          to trying to guess whatever number Bob has from $1$ through $x$ with
          constant cost and around $\frac{1}{\sqrt{x}}$ accuracy.
          Random guessing gets us constant cost, but only $\frac{1}{x}$ accuracy
          which extends to a protocol that's only correct\
          $\frac{1}{2}+\frac{1}{\Theta\left(2^n\right)}$ of the time.

          I also tried seeing if the amount of guesses needed to get at least
          one correct was bounded above, but unfortunately we need at least
          \[t \ge \frac{\ln \left(1-\frac{1}{2^{n/2}}\right)}{-\frac{1}{2^n}}
          = -2^n\ln\left(1-\frac{1}{2^{n/2}}\right)\]
          trials for the target correctness probability, which diverges as $n \to \infty$.

          Since each bit transmitted can remove at most half the possibilities,
          this seems impossible.
\end{enumerate}

\newpage

\section*{Problem 9}

WLOG let $F=\mathbb{R}$.
We proved in class that ranks are maximized when we consider them over $\mathbb{R}$,
so this makes the constraint on $c$ as tight as possible.

We earlier derived that $\rk_{\mathbb{F}_2} M_{IP_n}=n$, while $\rk_{\mathbb{R}} M_{IP_n}=2^n-1$.
This means that $2^n-1 \le c^n$ and $c \ge 2$ by extension, since any base less than $2$
would eventually get dominated as $n \to \infty$.

\newpage

\section*{Problem 10}

I implemented the LP using \verb|cvxpy| and saw that optimal solutions for $n=3$ and $n=4$
both had $\mu$ be the uniform distribution over all the inputs.
By engineer's induction, let's just say that it's true for all $n$.

Assuming that the uniform distribution is optimal, the rectangle that captures
the largest amount of it would be any $A \times B$ where $|A|=|B|$ and $B=A^C$,
since the size of a maximal rectangle is quadratic in terms of $|A|$.
This captures exactly $\frac{2^n}{2} \cdot \frac{2^n}{2}=2^{2n-2}$ cells.

Since each cell has the same probability of $\frac{1}{4^n-2^n}$, we have
\[RS(\lnot EQ_n)=\frac{2^{2n-2}}{4^n-2^n}=\boxed{\frac{2^{n-2}}{2^n-1}}\]
As $n \to \infty$, this expression approaches $\frac{1}{4}$, the lower
bound we derived in class, so it seems to make sense.

\end{document}
